\documentclass[UTF8,11pt]{ctexart}
\usepackage[a4paper,margin=24mm]{geometry}
\usepackage{amsmath,amssymb,amsthm,mathtools,mathrsfs}
\usepackage[all]{xy}
\usepackage{xurl,hyperref,enumitem}
\usepackage{tikz}
\usetikzlibrary{arrows.meta,cd,decorations.markings,calc,patterns}
\hypersetup{hidelinks,breaklinks=true}
\setlength{\emergencystretch}{3em}
\allowdisplaybreaks
\newtheorem*{theorem}{Theorem}
\newtheorem*{thm}{Theorem}
\newtheorem*{lemma}{Lemma}
\newtheorem*{proposition}{Proposition}
\newtheorem*{prop}{Proposition}
\newtheorem*{corollary}{Corollary}
\newtheorem*{cor}{Corollary}
\newtheorem*{claim}{Claim}
\theoremstyle{definition}
\newtheorem*{definition}{Definition}
\newtheorem*{defn}{Definition}
\newtheorem*{remark}{Remark}
\newtheorem*{example}{Example}
\newcommand{\R}{\mathbb R}
\newcommand{\C}{\mathbb C}
\newcommand{\Z}{\mathbb Z}
\newcommand{\Q}{\mathbb Q}
\newcommand{\N}{\mathbb N}
\newcommand{\F}{\mathbb F}
\newcommand{\D}{\mathbb D}
\newcommand{\bB}{\mathbb B}
\newcommand{\bC}{\mathbb C}
\newcommand{\bR}{\mathbb R}
\newcommand{\bZ}{\mathbb Z}
\newcommand{\bN}{\mathbb N}
\newcommand{\bQ}{\mathbb Q}
\newcommand{\bD}{\mathbb D}
\newcommand{\bF}{\mathbb F}
\newcommand{\CP}{\mathbb{CP}}
\newcommand{\RP}{\mathbb{RP}}
\newcommand{\sO}{\mathscr O}
\newcommand{\sI}{\mathscr I}
\newcommand{\sF}{\mathscr F}
\newcommand{\sG}{\mathscr G}
\newcommand{\sH}{\mathscr H}
\newcommand{\sR}{\mathscr R}
\newcommand{\sM}{\mathscr M}
\newcommand{\sV}{\mathscr V}
\newcommand{\sU}{\mathscr U}
\DeclareMathOperator{\Hom}{Hom}
\DeclareMathOperator{\Ext}{Ext}
\DeclareMathOperator{\Tor}{Tor}
\DeclareMathOperator{\Spec}{Spec}
\DeclareMathOperator{\Proj}{Proj}
\DeclareMathOperator{\supp}{supp}
\DeclareMathOperator{\im}{im}
\DeclareMathOperator{\id}{id}
\DeclareMathOperator{\Tr}{Tr}
\DeclareMathOperator{\tr}{tr}
\DeclareMathOperator{\rank}{rank}
\DeclareMathOperator{\ord}{ord}
\DeclareMathOperator{\dist}{dist}
\DeclareMathOperator{\End}{End}
\DeclareMathOperator{\Aut}{Aut}
\DeclareMathOperator{\Gal}{Gal}
\DeclareMathOperator{\vol}{vol}
\DeclareMathOperator{\Cl}{Cl}
\DeclareMathOperator{\coker}{coker}
\DeclareMathOperator{\codim}{codim}
\DeclareMathOperator{\divergence}{div}
\DeclareMathOperator{\Ric}{Ric}
\DeclareMathOperator{\grad}{grad}
\DeclareMathOperator{\Hess}{Hess}
\newcommand{\abs}[1]{\left\lvert#1\right\rvert}
\newcommand{\sabs}[1]{\lvert#1\rvert}
\newcommand{\babs}[1]{\bigl\lvert#1\bigr\rvert}
\newcommand{\Babs}[1]{\Bigl\lvert#1\Bigr\rvert}
\newcommand{\bbabs}[1]{\biggl\lvert#1\biggr\rvert}
\newcommand{\BBabs}[1]{\Biggl\lvert#1\Biggr\rvert}
\newcommand{\norm}[1]{\left\lVert#1\right\rVert}
\newcommand{\snorm}[1]{\lVert#1\rVert}
\newcommand{\bnorm}[1]{\bigl\lVert#1\bigr\rVert}
\newcommand{\Bnorm}[1]{\Bigl\lVert#1\Bigr\rVert}
\newcommand{\bbnorm}[1]{\biggl\lVert#1\biggr\rVert}
\newcommand{\BBnorm}[1]{\Biggl\lVert#1\Biggr\rVert}
\newcommand{\widebar}[1]{\overline{#1}}
\newcommand{\nicefrac}[2]{\frac{#1}{#2}}
\newcommand{\linnprod}[2]{\langle#1,#2\rangle}
\newcommand{\myindex}[1]{#1}
\newcommand{\glsadd}[1]{}
\newcommand{\avoidbreak}{}
\newcommand{\sourceref}[1]{\textnormal{[原文：\nolinkurl{#1}]}}
\newcommand{\thmref}[1]{\sourceref{#1}}
\newcommand{\lemmaref}[1]{\sourceref{#1}}
\newcommand{\propref}[1]{\sourceref{#1}}
\newcommand{\corref}[1]{\sourceref{#1}}
\newcommand{\defnref}[1]{\sourceref{#1}}
\newcommand{\exampleref}[1]{\sourceref{#1}}
\newcommand{\exerciseref}[1]{\sourceref{#1}}
\newcommand{\figureref}[1]{\sourceref{#1}}
\newcommand{\sectionref}[1]{\sourceref{#1}}
\newcommand{\appendixref}[1]{\sourceref{#1}}
\newcommand{\stacksref}[2]{\href{https://stacks.math.columbia.edu/tag/#1}{#2 [#1]}}

\begin{document}
\section*{050\quad 阿贝尔正规因子的半直积表示分类}
\subsection*{来源与计数目标}
Pavel Etingof 与 Oleg Golberg、Sebastian Hensel、Tiankai Liu、Alex Schwendner、Elena Yudovina、Dmitry Vaintrob 合著的 \emph{Introduction to Representation Theory} 课程讲义。MIT 18.712，Fall 2010 的 OCW 原讲义版本，不是限制转载的 2016 年 AMS 出版书。Chapter 4，\S4.26，Theorem 4.75(i)--(iii)，含其之前的诱导模型定义和对应证明，分章 PDF 第30--31页。。获取日期：2026-09-28。
\par\url{https://ocw.mit.edu/courses/18-712-introduction-to-representation-theory-fall-2010/84358595a02a73bced2c4e363a5d66f0_MIT18_712F10_ch4.pdf}
\par\url{https://ocw.mit.edu/courses/18-712-introduction-to-representation-theory-fall-2010/pages/lecture-notes/}
\par\textbf{本文件只计一个问题：}有限群 $G,A$、$A$ 阿贝尔时，以 $G$ 在 $A^*$ 上的轨道及稳定子不可约表示，构造并分类 $G\ltimes A$ 的全部复不可约表示。
\subsection*{作者命题与人类证明}

Let $G,A$ be groups and $\rho:G\to\Aut(A)$ be a homomorphism. For $a\in A$, denote $\rho(g)a$ by $g(a)$. The semidirect product $G\ltimes A$ is defined to be the product $A\times G$ with multiplication law
\[
(a_1,g_1)(a_2,g_2)=(a_1g_1(a_2),g_1g_2).
\]
Clearly, $G$ and $A$ are subgroups of $G\ltimes A$ in a natural way.

We would like to study irreducible complex representations of $G\ltimes A$. For simplicity, let us do it when $A$ is abelian. In this case, irreducible representations of $A$ are $1$-dimensional and form the character group $A^*$, which carries an action of $G$. Let $O$ be an orbit of this action, $x\in O$ a chosen element, and $G_x$ the stabilizer of $x$ in $G$. Let $U$ be an irreducible representation of $G_x$. Then we define a representation $V_{(O,U)}$ of $G\ltimes A$ as follows.

As a representation of $G$, we set
\[
V_{(O,x,U)}=\operatorname{Ind}_{G_x}^G U
=\{f:G\to U\mid f(hg)=hf(g),\ h\in G_x\}.
\]
Next, we introduce an additional action of $A$ on this space by $(af)(g)=x(g(a))f(g)$. Then it's easy to check that these two actions combine into an action of $G\ltimes A$. Also, it is clear that this representation does not really depend on the choice of $x$, in the following sense. Let $x,y\in O$, and $g\in G$ be such that $gx=y$, and let $g(U)$ be the representation of $G_y$ obtained from the representation $U$ of $G_x$ by the action of $g$. Then $V_{(O,x,U)}$ is (naturally) isomorphic to $V_{(O,y,g(U))}$. Thus we will denote $V_{(O,x,U)}$ by $V_{(O,U)}$ (remembering, however, that $x$ has been fixed).
\begin{theorem}[4.75, parts (i)--(iii)]
\begin{enumerate}[label=(\roman*)]
\item The representations $V_{(O,U)}$ are irreducible.
\item They are pairwise nonisomorphic.
\item They form a complete set of irreducible representations of $G\ltimes A$.
\end{enumerate}
\end{theorem}
\begin{proof}
(i) Let us decompose $V=V_{(O,U)}$ as an $A$-module. Then we get
\[
V=\bigoplus_{y\in O}V_y,
\]
where $V_y=\{v\in V_{(O,U)}\mid av=(y,a)v,\ a\in A\}$. (Equivalently, $V_y=\{v\in V_{(O,U)}\mid v(g)=0\text{ unless }gy=x\}$.) So if $W\subset V$ is a subrepresentation, then $W=\bigoplus_{y\in O}W_y$, where $W_y\subset V_y$. Now, $V_y$ is a representation of $G_y$, which goes to $U$ under any isomorphism $G_y\to G_x$ determined by $g\in G$ mapping $x$ to $y$. Hence, $V_y$ is irreducible over $G_y$, so $W_y=0$ or $W_y=V_y$ for each $y$. Also, if $hy=z$ then $hW_y=W_z$, so either $W_y=0$ for all $y$ or $W_y=V_y$ for all $y$, as desired.

(ii) The orbit $O$ is determined by the $A$-module structure of $V$, and the representation $U$ by the structure of $V_x$ as a $G_x$-module.

(iii) We have
\[
\begin{aligned}
\sum_{U,O}\dim V_{(U,O)}^2
&=\sum_{U,O}|O|^2(\dim U)^2
 =\sum_O|O|^2|G_x|\\
&=\sum_O|O||G/G_x||G_x|
 =|G|\sum_O|O|=|G||A^*|=|G\ltimes A|.
\end{aligned}
\]
\end{proof}

\subsection*{整理说明（非作者正文）}
有限性是本章有限群表示语境，原段开头只写 groups；在计数目标中显式恢复这一适用范围。这里截取的是作者定理的 (i)--(iii) 与其完整证明；(iv) 的特征标公式只写与 Mackey 公式证明相同，故不纳入本条目标，不将该部分冒称已展开。保留作者在同一段中两种下标次序 $V_{(O,U)}$、$V_{(U,O)}$ 的写法。

标准先修：有限阿贝尔群的复表示分解、诱导表示的函数模型、轨道--稳定子公式及不可约表示维数平方和等于群阶。实际独立工作是把群作用的轨道组织成表示的权空间，构造半直积作用、恢复稳定子表示，并用维数总量证明穷尽性，不是直接调用分类结论。正文文本已实际取得；两页截图未成功，未声称逐页图像检查。
\subsection*{可修改的简短标注（非作者正文）}
$c$：有限半直积群的不可约表示分类。

$a$：阿贝尔群特征标；权空间分解；群作用轨道与稳定子；诱导表示函数模型；维数平方和。

\texttt{cross\_theory\_status = yes}。将半直积表示限制到阿贝尔正规子群得到特征标权空间，再把权的轨道与稳定子表示对应回不可约模；轨道计数给出穷尽性。位置：4.26 的模型构造及 4.75(i)--(iii)。

全部标注均为非穷尽、可修改初稿；未标注不表示未使用。
\subsection*{许可与修改记录}
材料来自 MIT OpenCourseWare，作者与课程见来源。按该站所示 CC BY-NC-SA 4.0 许可复用，整理部分沿用同一许可。\url{https://creativecommons.org/licenses/by-nc-sa/4.0/}。2026-09-28：助手转录题证、调整数学排版并加入来源、范围和初稿标注；不暗示作者或 MIT 认可本次整理。所留为本次题证转录摘录，不是原 PDF 下载件。
\end{document}
